\bookchapter{Installing Software and Working Remotely}{ch:07-installing-software-and-working-remotely}

\begin{chapterintro}

This chapter covers two everyday jobs: getting the software you need, safely and in the right version, and working on a machine that isn't in front of you, with SSH keys, file copies, and sessions that survive a dropped connection.

\end{chapterintro}

\emph{The problem.} I need a tool on a machine I can only reach over the internet.

\emph{The question.} How do I install tools safely, and work on a machine that isn't in front of me?

\section{Package Managers}\label{07-installing-software-and-working-remotely:package-managers}

On a phone you use the app store, not random websites. The terminal's app store is a \textbf{package manager}\index{package manager}. It downloads software from a trusted source, installs it, and can update or remove it later. Each piece of software is a package, and most need other packages, their dependencies, which the package manager installs for you.

\Needspace{9\baselineskip}

\begin{booktable}
\begin{longtable}{>{\raggedright\arraybackslash}p{\dimexpr 0.2212\linewidth-1.5577\tabcolsep\relax}>{\raggedright\arraybackslash}p{\dimexpr 0.1827\linewidth-1.6346\tabcolsep\relax}>{\raggedright\arraybackslash}p{\dimexpr 0.2212\linewidth-1.5577\tabcolsep\relax}>{\raggedright\arraybackslash}p{\dimexpr 0.3750\linewidth-1.2500\tabcolsep\relax}}
\tabletop
\rowcolor{tablehead}\thead{System} & \thead{Package manager} & \thead{Install} & \thead{Update everything} \\
\tablemid
\endfirsthead
\tabletop
\rowcolor{tablehead}\thead{System} & \thead{Package manager} & \thead{Install} & \thead{Update everything} \\
\tablemid
\endhead
\tablebottom
\endlastfoot
Ubuntu, Debian, WSL & \texttt{apt} & \texttt{sudo\ apt\ install\ jq} & \texttt{sudo\ apt\ update\ \&}\allowbreak{}\texttt{\&}\allowbreak{}\texttt{\ sudo\ apt\ upgrade} \\
Fedora, Red Hat & \texttt{dnf} & \texttt{sudo\ dnf\ install\ jq} & \texttt{sudo\ dnf\ upgrade} \\
macOS & Homebrew & \texttt{brew\ install\ jq} & \texttt{brew\ update\ \&}\allowbreak{}\texttt{\&}\allowbreak{}\texttt{\ brew\ upgrade} \\
\end{longtable}
\end{booktable}

On Debian and Ubuntu, \texttt{apt\ update} only refreshes the list of what's available; \texttt{apt\ upgrade} installs newer versions. Run \texttt{update} before installing. Homebrew, installed from brew.sh, keeps everything in your own folders, so it needs no \texttt{sudo}.

\begin{callout}{warning}

Some guides say to paste \texttt{curl\ ...\ \textbar{}\ sudo\ bash}, which runs a downloaded script as root, unread. Prefer your package manager, or download the script and read it first.

\end{callout}

\section{Virtual Environments for Python}\label{07-installing-software-and-working-remotely:virtual-environments-for-python}

Project A needs version 1 of a library and project B needs version 2, but your system has one Python. A \textbf{virtual environment}\index{virtual environment} solves this: a folder with its own Python and its own packages, one per project. Create one in the project folder, then activate it:

\codeneed{5}\begin{codeblock}

\begin{Shaded}
\begin{Highlighting}[]
\ExtensionTok{python3} \AttributeTok{{-}m}\NormalTok{ venv .venv}
\BuiltInTok{source}\NormalTok{ .venv/bin/activate}
\FunctionTok{which}\NormalTok{ python}
\end{Highlighting}
\end{Shaded}

\end{codeblock}

\codeneed{1}\begin{codeblock}
\begin{Verbatim}[formatcom=\outputstyle]
/home/ada/sandbox/.venv/bin/python
\end{Verbatim}
\end{codeblock}

Activating puts \texttt{.venv/bin} at the front of your PATH, so \texttt{python} and \texttt{pip} mean the project's copies, and your prompt shows \texttt{(.venv)}. Now \texttt{pip\ install\ requests} installs into this project only. \texttt{deactivate} goes back. A version manager, such as \texttt{pyenv} or \texttt{uv} for Python, goes further and installs whole language versions side by side.

\begin{callout}{tip}

Keep \texttt{.venv} out of Git (Chapter 8 shows how). Record what the project needs with \texttt{pip\ freeze\ \textgreater{}\ requirements.}\allowbreak{}\texttt{txt} instead, so anyone can rebuild the environment.

\end{callout}

\section{SSH: A Terminal on Another Machine}\label{07-installing-software-and-working-remotely:ssh-a-terminal-on-another-machine}

\textbf{SSH}\index{ssh@SSH} (Secure Shell) lets you open a shell on another computer over the network. Everything you type and see travels encrypted.

\codeneed{1}\begin{codeblock}

\begin{Shaded}
\begin{Highlighting}[]
\FunctionTok{ssh}\NormalTok{ ada@203.0.113.10}
\end{Highlighting}
\end{Shaded}

\end{codeblock}

That means ``log in as \texttt{ada} on the machine at that address''. The first time, SSH shows the server's fingerprint and asks whether to trust it; it remembers your answer in \texttt{\textasciitilde{}/}\allowbreak{}\texttt{.ssh/}\allowbreak{}\texttt{known\_hosts} and warns loudly if the fingerprint ever changes.

You could log in with a password, but passwords can be guessed. The better way is a \textbf{key pair}\index{key pair}: two linked files. Think of a padlock and its key. The \textbf{public key}\index{public key} is the padlock: you can hand out copies freely and fit one to every server you use. The \textbf{private key}\index{private key} is the only key that opens those padlocks, and it never leaves your laptop.

\codeneed{1}\begin{codeblock}

\begin{Shaded}
\begin{Highlighting}[]
\FunctionTok{ssh{-}keygen} \AttributeTok{{-}t}\NormalTok{ ed25519 }\AttributeTok{{-}C} \StringTok{"ada@laptop"} \AttributeTok{{-}f}\NormalTok{ \textasciitilde{}/.ssh/id\_ed25519}
\end{Highlighting}
\end{Shaded}

\end{codeblock}

\codeneed{6}

It asks for a passphrase twice. Choose one: it protects the private key if your laptop is stolen. Then:

\codeneed{3}\begin{codeblock}
\begin{Verbatim}[formatcom=\outputstyle]
Your identification has been saved in /home/ada/.ssh/id_ed25519
Your public key has been saved in /home/ada/.ssh/id_ed25519.pub
…
\end{Verbatim}
\end{codeblock}

\codeneed{6}

The \texttt{.pub} file is the padlock:

\codeneed{3}\begin{codeblock}

\begin{Shaded}
\begin{Highlighting}[]
\FunctionTok{cat}\NormalTok{ \textasciitilde{}/.ssh/id\_ed25519.pub}
\end{Highlighting}
\end{Shaded}

\end{codeblock}

\codeneed{1}\begin{codeblock}
\begin{Verbatim}[formatcom=\outputstyle]
ssh-ed25519 AAAAC3NzaC1lZDI1NTE5AAAAIDIpXo4KbS5P4qftLjgXDY393xIkLDJdG/ggTPq9Dqyk ada@laptop
\end{Verbatim}
\end{codeblock}

\texttt{ssh-}\allowbreak{}\texttt{copy-}\allowbreak{}\texttt{id\ ada@203.}\allowbreak{}\texttt{0.}\allowbreak{}\texttt{113.}\allowbreak{}\texttt{10} installs it on a server; GitHub and cloud providers have a page where you paste it. From then on, \texttt{ssh} logs you in with the key.

\begin{callout}{warning}

Never share or copy \texttt{\textasciitilde{}/.ssh/id\_ed25519} (the file without \texttt{.pub}). Anyone who has it can log in as you everywhere you've installed the padlock. \texttt{ssh-keygen} makes it \texttt{600} for you; keep it that way.

\end{callout}

\codeneed{6}

To save typing, name your servers in \texttt{\textasciitilde{}/.ssh/config}:

\codeneed{3}\begin{codeblock}
\begin{Verbatim}[formatcom=\outputstyle]
Host web
    HostName 203.0.113.10
    User ada
\end{Verbatim}
\end{codeblock}

Now \texttt{ssh\ web} is enough.

\codeneed{7}

\section{Copying Files Between Machines}\label{07-installing-software-and-working-remotely:copying-files-between-machines}

\texttt{scp} copies files over SSH, like \texttt{cp} with a remote path written \texttt{host:path}: \texttt{scp\ report.}\allowbreak{}\texttt{csv\ web:}\allowbreak{}\texttt{/}\allowbreak{}\texttt{home/}\allowbreak{}\texttt{ada/}. For folders you copy again and again, use rsync, which sends only what changed. It works between local folders too:

\codeneed{1}\begin{codeblock}

\begin{Shaded}
\begin{Highlighting}[]
\FunctionTok{rsync} \AttributeTok{{-}av}\NormalTok{ notes/ backup/notes/}
\end{Highlighting}
\end{Shaded}

\end{codeblock}

It lists each file it copies. Add a file to \texttt{notes} and run the same command again, and only the new file is sent. The trailing \texttt{/} on \texttt{notes/} means ``the contents of notes''.

\section{Keeping Work Alive with tmux}\label{07-installing-software-and-working-remotely:keeping-work-alive-with-tmux}

If your connection drops, the shell on the server ends, and so does anything running in its foreground. \texttt{tmux} runs sessions on the server that keep going while you're away. Start one with \texttt{tmux\ new\ -s\ work}, run your long job, then press Ctrl-B followed by D to detach. Later, from a new login, \texttt{tmux\ attach\ -t\ work} puts you back exactly where you were. Install it with \texttt{apt\ install\ tmux} or \texttt{brew\ install\ tmux}.

\section{Lab: Your Toolbox}\label{07-installing-software-and-working-remotely:lab-your-toolbox}

Work in \texttt{\textasciitilde{}/sandbox}. None of these steps needs a server.

\begin{enumerate}
\def\labelenumi{\arabic{enumi}.}
\tightlist
\item
  Install jq and shellcheck with your package manager if you haven't: \texttt{sudo\ apt\ install\ jq\ shellcheck} or \texttt{brew\ install\ jq\ shellcheck}. \textbf{Expected:} \texttt{jq\ -\/-version} prints a version such as \texttt{jq-1.8.2}.
\item
  Create and activate a virtual environment: \texttt{python3\ -m\ venv\ .venv} then \texttt{source\ .}\allowbreak{}\texttt{venv/}\allowbreak{}\texttt{bin/}\allowbreak{}\texttt{activate}. \textbf{Expected:} your prompt starts with \texttt{(.venv)}, and \texttt{which\ python} points inside \texttt{\textasciitilde{}/sandbox/.venv}.
\item
  Run \texttt{pip\ list}, then \texttt{deactivate}. \textbf{Expected:} only \texttt{pip} is listed, then \texttt{(.venv)} disappears.
\item
  Make a practice key pair: \texttt{ssh-}\allowbreak{}\texttt{keygen\ -}\allowbreak{}\texttt{t\ ed25519\ -}\allowbreak{}\texttt{C\ "practice"\ -}\allowbreak{}\texttt{f\ \textasciitilde{}/}\allowbreak{}\texttt{sandbox/}\allowbreak{}\texttt{practice\_}\allowbreak{}\texttt{key}. \textbf{Expected:} two files, \texttt{practice\_key} (mode \texttt{-rw-\/-\/-\/-\/-\/-\/-}) and \texttt{practice\_key.pub}. Check with \texttt{ls\ -l\ practice\_key*}.
\item
  Read the public half: \texttt{cat\ practice\_key.pub}. \textbf{Expected:} one line starting with \texttt{ssh-ed25519}. Delete both files when you're done.
\item
  Sync a folder twice: \texttt{mkdir\ -p\ notes}, \texttt{echo\ a\ \textgreater{}\ notes/}\allowbreak{}\texttt{monday.}\allowbreak{}\texttt{txt}, \texttt{rsync\ -}\allowbreak{}\texttt{av\ notes/}\allowbreak{}\texttt{\ backup/}\allowbreak{}\texttt{notes/}; then \texttt{echo\ c\ \textgreater{}\ notes/}\allowbreak{}\texttt{wednesday.}\allowbreak{}\texttt{txt} and run the same \texttt{rsync} again. \textbf{Expected:} the second run copies only \texttt{wednesday.txt} (it may also list \texttt{./}, the folder itself).
\end{enumerate}

\section{Summary, Key Terms, and Review Questions}\label{07-installing-software-and-working-remotely:summary-key-terms-and-review-questions}

\subsection*{Summary}\label{07-installing-software-and-working-remotely:summary}

\begin{itemize}
\tightlist
\item
  Package managers install software and its dependencies from trusted sources: \texttt{apt}, \texttt{dnf}, Homebrew.
\item
  A virtual environment gives each Python project its own packages; version managers install whole language versions.
\item
  SSH opens an encrypted shell on another machine. Key pairs beat passwords: share the public key, guard the private key.
\item
  \texttt{scp} copies files over SSH; \texttt{rsync} copies only what changed.
\item
  \texttt{tmux} keeps sessions alive when your connection drops.
\end{itemize}

\Needspace{13\baselineskip}

\subsection*{Key Terms}\label{07-installing-software-and-working-remotely:key-terms}

\begin{booktable}
\begin{longtable}{>{\raggedright\arraybackslash}p{\dimexpr 0.2644\linewidth-1.4713\tabcolsep\relax}>{\raggedright\arraybackslash}p{\dimexpr 0.7356\linewidth-0.5287\tabcolsep\relax}}
\tabletop
\rowcolor{tablehead}\thead{Term} & \thead{Meaning} \\
\tablemid
\endfirsthead
\tabletop
\rowcolor{tablehead}\thead{Term} & \thead{Meaning} \\
\tablemid
\endhead
\tablebottom
\endlastfoot
Package manager & A tool that installs, updates, and removes software \\
Virtual environment & A per-project folder with its own Python and packages \\
SSH & Secure Shell: an encrypted way to use another machine's shell \\
Key pair & A linked public and private key used to prove who you are \\
Public key & The half you share, installed on servers like a padlock \\
Private key & The secret half that stays on your machine \\
\end{longtable}
\end{booktable}

\subsection*{Review Questions}\label{07-installing-software-and-working-remotely:review-questions}

\begin{enumerate}
\def\labelenumi{\arabic{enumi}.}
\tightlist
\item
  What's the difference between \texttt{apt\ update} and \texttt{apt\ upgrade}?
\item
  Why does each Python project deserve its own virtual environment?
\item
  In the padlock picture, which file do you give to a server, and which do you never share?
\item
  When would you choose \texttt{rsync} over \texttt{scp}?
\item
  You start a two-hour job over SSH and your Wi-Fi drops. How could \texttt{tmux} have saved it?
\end{enumerate}

\begin{understandbox}

\textbf{You understand this when} you can set up a project's Python environment, and explain how an SSH key lets you in without a password.

\end{understandbox}
